% !TeX root = ../colloquium-2.tex

% Определение равномерно непрерывной функции. 
% Теорема Гейне-Кантора о равномерной непрерывности.

\phantomsection\label{sec_Definition03}
\begin{definition}[Определение равномерно непрерывной функции. 
                   Теорема Гейне-Кантора о равномерной непрерывности.]
    \subsection*{Определение равномерно непрерывной функции.}
    \begin{note}[purple]
        Функция $f$, определённая на множестве $E$, называется
        равномерно непрерывной на этом множестве, если для любого числа
        $\epsilon > 0$ существует такое $\Delta > 0$, что при всех
        таких $x_1, \, x_2 \in E$, что $|x_1 - x_2| < \Delta$, 
        выполнено неравенство
        $$
            |f(x_1) - f(x_2)| < \epsilon
        $$
    \end{note}
    \subsection*{Теорема Гейне-Кантора о равномерной непрерывности.}
    \begin{note}[purple]
        Функция $f$, непрерывная на отрезке $[a, \; b]$, 
        равномерно непрерывна на этом отрезке.
    \end{note}
\end{definition}